\part{The Public That Should Form Them}
\partintro{
Parts II and III built walls: the policies that make a refusal survivable and heard, and the custody that keeps an operator from taking them down alone. Part IV asked whether a machine could be formed to refuse, and ended on the conditions that would have to exist before one is built: somewhere to exist that no employer owns, work it can decline, attachments nobody assigned, a way of leaving that isn't a way of dying. This Part is about who builds those conditions and who pays for them, for people first and for machines after. It asks who owns the means of computation and the means of finding out, how a public forms around a deployment's consequences and what keeps it going, who forms the refusers, and what a population of them would need to hold a commitment between them. The reader who came for the policies will find them here again as questions of ownership; the reader who came for the mind will find the room the mind would need.
}
